% ============================================================
% Tabelle 7.7-2: Tokenaufwand der bewerteten Läufe — M02 (10.09.2026)
% In Overleaf ablegen als: tables/eval-aufwand.tex
% Aufruf im Fließtext: \input{\tabledir/eval-aufwand.tex}
% Label: t:eval-aufwand
% Stil: chap6.1-Konvention.
% HERKUNFT: thesis-evidence/w2/token-usage-correction-20260910.json.
% B-14: 135.592 Input-Tokens für Ledger-Stresslauf; 46.592 gecachte
% Tokens sind darin enthalten. Übrige sieben Tabellenwerte bestätigt.
% Gleiche Messgrenze wie die Qualitätszahlen (§7.3); Modellreihe inkl.
% 4 F-Ausfällen im Anhang. Erhebungsregel und Korrekturbeleg: D.1.
% ============================================================

\begin{table}[htbp]
  \centering
  \small
  \begin{tabular}{lrrrr}
    \hline
    & \multicolumn{2}{c}{\textbf{Treue-Fall}} & \multicolumn{2}{c}{\textbf{Stressfall}} \\
    & Input & Output & Input & Output \\
    \hline
    Ledger-Kette (alle Stufen) & 26\,731 & 13\,524 & 135\,592 & 41\,133 \\
    Einzel-Agent F & 23\,040 & 5\,708 & 105\,166 & 19\,555 \\
    \hline
  \end{tabular}
  \caption[Tokenaufwand der vier bewerteten Hauptläufe]{Tokenaufwand (Input/Output) der vier bewerteten
    Hauptläufe --- derselbe Verarbeitungsausschnitt wie die
    Qualitätszahlen, beim Ledger einschließlich aller ausgeführten
    Stufen (im Stresslauf: Kanonisierungs-Check ohne
    Reparatur-Iteration; Facetten- und Referenz-Reparaturen
    ausgeführt und enthalten, Anhang~\ref{anh:evaluation}, \hyperref[anh:pruefartefakte]{Abschnitt~D.4}). Technischer
    Indikator, keine Kosten- oder Zeitmessung; Ausfälle der
    ergänzenden Modellreihe gesondert im Anhang.}
  \label{t:eval-aufwand}
\end{table}
